\documentclass[11pt]{article}
\usepackage[margin=1in]{geometry}
\usepackage[T1]{fontenc}\usepackage[utf8]{inputenc}\usepackage{lmodern}
\usepackage{amsmath,amssymb,amsthm,mathtools,microtype}
\usepackage[colorlinks=true,linkcolor=blue,urlcolor=blue]{hyperref}
\hypersetup{pdftitle={Formula complexity and prior probabilities in Bayesian neural networks},pdfauthor={ChatGPT-6 Astra; Claude Opus 5.5; Kaarel H\"anni; Dmitry Vaintrob}}
\newtheorem{theorem}{Theorem}[section]\newtheorem{lemma}[theorem]{Lemma}\newtheorem{corollary}[theorem]{Corollary}\newtheorem{proposition}[theorem]{Proposition}
\newcommand{\E}{\mathbb E}\newcommand{\R}{\mathbb R}\newcommand{\KL}{\operatorname{KL}}
\title{Formula complexity and prior probabilities\\in Bayesian neural networks}
\author{ChatGPT-6 Astra \and Claude Opus 5.5 \and Kaarel H\"anni \and Dmitry Vaintrob\thanks{Authors are listed in alphabetical order.}}
\date{17 November 2026}
\begin{document}\maketitle
\begin{abstract}
We relate actual neural-network prior probabilities to finite-precision formula complexity in both directions. The network has one fixed dense architecture, and all weight coordinates are independent and absolutely continuous. A coordinate of fan-in $F$ is an increment over time $1/F$ of a fixed symmetric finite-variation L\'evy process. For any globally Lipschitz activation, a small bounded-coefficient formula has a uniform function neighborhood whose negative log prior mass is polynomial in formula size, depth and input dimension, independently of learner width. Conversely, a depth-$H$ prior draw is, with probability $1-\delta$, uniformly approximable to accuracy $\tau$ by a rational-weight formula of size at most $[CL\log(eL)]^{H+2}$ and coefficient bit length at most $CL^3$, where $L=H+1+\log(1/(\delta\tau))$. A target neighborhood of prior mass $p$ therefore has a short formula approximation with resources polynomial in $\log(1/p)$ at fixed depth. The converse follows by deleting tiny Poisson jumps, exploring only ancestors of the output, and quantizing retained weights. For ReLU and Boolean inputs, a margin version yields an ordinary Boolean-circuit consequence. An explicit signed beta-prime coordinate density satisfies the same conclusions. At fixed depth, the wide L\'evy prior converges to a random-formula law; a finite-width motif estimate explains why directly retaining a shared computation incurs an additional logarithmic width cost. No robustness restriction or prior over network architectures is used.
\end{abstract}

\section{The prior and the two directions}
A function prior concerns the probabilities assigned to functions, not only a learning bound. We study uniform neighborhoods of real-valued network functions. A small implementation gives a lower bound on neighborhood probability. Conversely, a neighborhood with appreciable prior probability contains a function with a small finite-precision implementation. At fixed depth these bounds give a polynomial relation between approximate formula complexity and negative log prior mass.

Fix a globally $A$-Lipschitz activation $\phi:\R\to\R$, with $A\ge1$ and $b=|\phi(0)|$. Inputs lie in $[-1,1]^d$. A width-$N$ network with $H$ hidden layers and a linear readout computes
\[
 h^1=\phi(W^1(1,x)),\quad h^\ell=\phi(W^\ell h^{\ell-1}),
 \quad z_w=v^{\mathsf T}h^H.
\]
Only the first layer has a constant input coordinate. Fix $c>0$ and $p>2$. Let $L_t$ be the absolutely convergent Poisson sum with symmetric jump intensity
\[
 t\,\mathfrak n(du),\qquad
 \mathfrak n(du)=\frac{c\,du}{|u|(1+|u|)^p}.
\]
A weight of fan-in $F$ has law $L_{1/F}$, independently of every other coordinate. Thus $F=d+1$ initially and $F=N$ later, including the readout. Write $\mu_N$ for this product prior. Every coordinate is nonzero almost surely. The jump representation describes the coordinate law; it does not choose a network architecture.

A layered formula duplicates any hidden computation used more than once, so every relevant hidden occurrence has one outgoing edge. Let $S\ge1$ count these occurrences, $m$ be its maximum hidden-layer width, and $B\ge1$ its coefficient bound. We assume $m\le N$. Unused vertices and zero edges may be deleted; zero-input hidden gates can remain when $\phi(0)\ne0$. Theorem~\ref{thm:uniform-mass} proves a bound of the form
\begin{align*}
 -\log\mu_N\{\|z_w-g\|_\infty\le\tau\}
 \le C(d+1)(S+1)\Bigg[&\log\frac{(S+1)(d+1)}\tau\\
 &+(H+2)\log\{C(1+A+b)(m+1)(B+8)\}\Bigg].
\end{align*}
For a shared width-$m$ circuit, use $S=\sum_{\ell=1}^H m^\ell$. This conversion can be exponential in depth; formula size is not identified with circuit size when computations are shared.

Theorem~\ref{thm:levy-prior-compression} gives the converse through a finite dictionary approximating almost every prior function. Its descriptions include input-coordinate addresses, and all bounds are independent of learner width. Corollary~\ref{cor:levy-prior-complexity-converse} applies the compression to any target neighborhood, including neighborhoods measured only on a prescribed subset such as the Boolean cube. Theorem~\ref{thm:margin-sign-converse} additionally treats the exact sign-function prior after conditioning on an all-input margin event; it needs no upper bound on raw outputs.

The activation may be bounded or unbounded, and need not vanish at zero. ReLU, GELU, tanh, sigmoid, softplus and SiLU satisfy the Lipschitz hypothesis. A fixed Lipschitz output map can be appended by adjusting tolerances. These results do not equate prior probability with an exponential of exact Boolean description length, and they do not remove depth dependence from the converse.

\section{One continuous infinitely divisible prior for every coordinate}
\label{sec:levy-scalar}
Fix constants $c>0$ and $p>2$, and define the symmetric L\'evy measure
\[
 \mathfrak n(dx)=\frac{c\,dx}{|x|(1+|x|)^p}
 \qquad(x\ne0).
\]
Let $L_t$ be the sum of all points of a Poisson random measure on $\R\setminus\{0\}$ with intensity $t\mathfrak n$, with no Gaussian component or deterministic drift. This sum converges absolutely almost surely. A coordinate whose fan-in is $F\ge1$ has the law
\begin{equation}
 \mu_F=\mathcal L(L_{1/F}).
 \label{eq:levy-coordinate-prior}
\end{equation}
All coordinates in the network are independent. Thus one fixed infinitely divisible family is used in every layer; fan-in changes its time parameter. The Poisson representation is a proof tool for this scalar law and does not introduce a prior over network architectures.

\begin{lemma}[Regularity and moments of the scalar prior]
\label{lem:levy-scalar-properties}
For every $t>0$, the law of $L_t$ is symmetric, centered, absolutely continuous with respect to Lebesgue measure, and has full support on $\R$. It has finite first and second absolute moments. If
\[
 M_j=\int |x|^j\mathfrak n(dx),\qquad j=1,2,
\]
then
\[
 \E|L_t|\le tM_1,\qquad \operatorname{Var}(L_t)=tM_2.
\]
For $0<r\le1$, write $\Lambda(r)=\mathfrak n\{|x|>r\}$. Constants depending only on $c,p$ satisfy
\begin{align*}
 \Lambda(r)&\le C[1+\log(1/r)],\\
 \int_{|x|\le r}|x|\mathfrak n(dx)&\le Cr,\\
 \int_{|x|\le r}x^2\mathfrak n(dx)&\le Cr^2.
\end{align*}
The density need not be bounded near zero, especially at small time.
\end{lemma}
\begin{proof}
Near zero the measure has density comparable to $1/|x|$, and at infinity it is comparable to $|x|^{-p-1}$. Hence $M_1,M_2$ are finite, and the three displayed estimates follow by direct integration. The expected total absolute jump size is $tM_1$, so the jump sum is absolutely convergent almost surely and its first moment has the stated bound. Symmetry and absolute integrability imply zero mean. For a finite jump cutoff the usual Poisson variance identity gives the second-moment formula; letting the cutoff decrease to zero in $L^2$ gives the formula for the full sum. Superposition of independent Poisson measures also proves infinite divisibility.

For fixed $r>0$, the jumps with magnitude above $r$ form a finite Poisson process. Conditional on at least one such jump, their sum has a density: conditional on its positive count, every jump has the absolutely continuous law $\mathfrak n(dx)\mathbf1_{|x|>r}/\Lambda(r)$. Adding the independent small-jump sum preserves absolute continuity. Therefore, for every Lebesgue-null set $A$,
\[
 \Pr(L_t\in A)\le e^{-t\Lambda(r)}.
\]
Since $\Lambda(r)\to\infty$ as $r\downarrow0$, the right side tends to zero.

To see full support, choose any nonzero $a$ in a prescribed open interval and a sufficiently small $\delta>0$ so that $[a-2\delta,a+2\delta]$ lies in the interval and excludes zero. Choose $r<\delta$ small enough that the expected absolute small-jump sum is less than $\delta/2$. Its magnitude is then at most $\delta$ with positive probability. Independently prescribe exactly one jump above $r$, lying in $(a-\delta,a+\delta)$, and no other jump above $r$. This has positive Poisson probability because the L\'evy density is positive away from zero. Their sum belongs to the prescribed interval with positive probability.
\end{proof}

\subsection{Suppressing a background coordinate}
\begin{lemma}[A small centered background law]
\label{lem:levy-background}
There is a constant $C$ depending only on $c,p$ such that, for every $F\ge1$ and $0<\epsilon\le1$, a law $Q_{F,\epsilon}^{\rm bg}\ll\mu_F$ satisfies
\begin{align*}
 &\E W=0,\qquad |W|\le\epsilon,\\
 &\E|W|\le C\epsilon/F,\qquad
   \E W^2\le C\epsilon^2/F,\\
 &\KL(Q_{F,\epsilon}^{\rm bg}\|\mu_F)
       \le\frac C F\log(C/\epsilon).
\end{align*}
\end{lemma}
\begin{proof}
Choose a fixed sufficiently large $\kappa$ depending only on $c,p$, and put $r=\epsilon/\kappa$. Split the Poisson sum into independent parts $S_r+J_r$, containing jumps of magnitude at most $r$ and above $r$, respectively. Force the large-jump process to have no points, and condition the small-jump sum on $|S_r|\le\epsilon$. Its mean remains zero by symmetry. Chebyshev and Lemma~\ref{lem:levy-scalar-properties} give
\[
 \Pr(|S_r|>\epsilon)
 \le\frac C{F\kappa^2}\le\frac1{4F}.
\]
The conditioning probability is at least $3/4$. The expected absolute small-jump sum is at most $Cr/F$, and its variance is at most $Cr^2/F$; division by the conditioning probability proves the first- and second-moment bounds.

On the latent Poisson space, the relative entropy of this joint conditioning is exactly
\[
 \frac{\Lambda(r)}F-
            \log\Pr(|S_r|\le\epsilon).
\]
The second term is at most $C/F$ because the failure probability is at most $1/(4F)$. The first term is at most $C\log(C/\epsilon)/F$. Mapping a Poisson realization to its summed weight cannot increase relative entropy, proving the stated scalar-law bound.
\end{proof}

\subsection{Placing a designated coefficient}
\begin{lemma}[An unbiased designated-coordinate law]
\label{lem:levy-designated}
For every $F\ge1$, $B\ge1$, $|a|\le B$, and $0<\epsilon\le1$, there is a law $Q_{F,a,\epsilon}^{\rm act}\ll\mu_F$ with
\begin{align*}
 &\E W=a,\qquad \epsilon<|W|\le B+8,\\
 &\E(W-a)^2\le C\epsilon^2,\\
 &\KL(Q_{F,a,\epsilon}^{\rm act}\|\mu_F)
       \le\log F+C\log\{C(B+8)/\epsilon\}.
\end{align*}
The construction includes $a=0$. In particular its support is disjoint from the background law's support.
\end{lemma}
\begin{proof}
Choose a designated jump $J$ with density $q$ as follows. If $|a|>4\epsilon$, take it uniform on $[a-\epsilon/2,a+\epsilon/2]$. Otherwise take the mixture of the uniform laws on $[5\epsilon,7\epsilon]$ and $[-7\epsilon,-5\epsilon]$, with positive-interval probability
\[
 \frac12\left(1+\frac a{6\epsilon}\right).
\]
This probability belongs to $[1/6,5/6]$. Both constructions have
\[
 \E J=a,\qquad \|q\|_\infty\le1/\epsilon,
 \qquad |J|>3\epsilon,\qquad |J|\le B+7,
 \qquad \operatorname{Var}(J)\le C\epsilon^2.
\]
Independently retain the original small-jump process below $r=\epsilon/\kappa$, conditioned on its sum $R$ satisfying $|R|\le\epsilon/2$. Enlarge the fixed $\kappa$ so that this event has probability at least $1-1/(4F)$. Its mean is zero and its variance is at most $C\epsilon^2/F$. Set $W=J+R$. This gives the asserted mean and variance. In the first case $|J|>7\epsilon/2$, and in the second $|J|\ge5\epsilon$; subtracting $|R|\le\epsilon/2$ gives the strict support separation. The upper support bound is immediate.

The entropy estimate is most transparent on the latent Poisson space. The original law gives an unordered configuration consisting of exactly one jump above $r$, at position $j$, the density
\[
 e^{-\Lambda(r)/F}\frac{\mathfrak n(dj)}F.
\]
The comparison instead prescribes that single jump with density $q(j)\,dj$, and independently conditions the original small-jump process as above. Its exact relative entropy is therefore
\begin{equation}
 \int q(j)\log\frac{Fq(j)}{\mathfrak n(j)}\,dj
       +\frac{\Lambda(r)}F
       -\log\Pr(|S_r|\le\epsilon/2),
 \label{eq:levy-latent-active-entropy}
\end{equation}
where $\mathfrak n(j)$ denotes the L\'evy density. There is no additional factorial term because the prescribed large-jump count is one. On the support of $q$,
\[
 \mathfrak n(j)\ge\frac c{(B+8)(B+9)^p}.
\]
Together with $q\le1/\epsilon$, this bounds the integral by
$\log F+\log(1/\epsilon)+C+C\log(B+8)$.
The two remaining terms are bounded exactly as in the background proof. Finally apply data processing to the summed scalar weight. This proves the result without requiring any lower bound on the density of $L_{1/F}$ itself.
\end{proof}

\section{Embedding a formula without a width cost}
\label{sec:embedding}
Fix a formula comparator and a precision $0<\epsilon\le1$. For each teacher coefficient $a$, use the designated-coordinate law from the preceding section. For a background coordinate, use the centered background law with cutoff $\epsilon$. Every designated weight has absolute value greater than $\epsilon$, while every background weight has absolute value at most $\epsilon$.

First suppose $N\ge2(S+1)$. Partition the indices in each hidden layer into disjoint groups assigned to that layer's formula occurrences. Each group has cardinality at least
$\lfloor N/(S+1)\rfloor\ge N/[2(S+1)]$. Independently choose one index uniformly from each group. These choices define an embedding $\iota$.

At every selected hidden row, use designated laws for the teacher's incoming edges and background laws for all other coordinates. A selected first-layer row uses designated laws for its $d+1$ input coefficients. Use designated laws at the readout for the teacher's final edges, and background laws on its remaining coordinates. Leave all unselected hidden rows at the native prior. For a fixed embedding, this comparison measure, denoted $Q_\iota$, is a product law on the weights. Its coordinates do not depend on the data or on preceding features.

\begin{lemma}[Disjoint supports and embedding entropy]\label{lem:embedding}
Distinct embeddings have disjoint supports under their comparison laws. Consequently, if $Q$ is their uniform mixture,
\begin{equation}
 \KL(Q\|\mu_N)
 =\frac1{|\mathcal I|}\sum_{\iota\in\mathcal I}
       \KL(Q_\iota\|\mu_N)-\log|\mathcal I|.
 \label{eq:embedding-entropy}
\end{equation}
Every designated internal or readout edge's $\log N$ cost is canceled by the entropy of its source occurrence's selected location, up to an additive $\log(2(S+1))$.
\end{lemma}
\begin{proof}
Start at the readout. In each group corresponding to one of its formula children, exactly one incoming coordinate has absolute value greater than $\epsilon$; that coordinate identifies the selected child. Inspect that row to identify its children in the same way. Every relevant hidden formula occurrence has one parent, and every relevant occurrence is reached from the output. Thus the weights determine the complete embedding. This proves disjointness.

On the support of $Q_\iota$, the mixture density is its density divided by $|\mathcal I|$. Integrating its logarithm gives~\eqref{eq:embedding-entropy}. There is exactly one outgoing edge for every hidden occurrence. Hence there are $S$ designated edges entering later hidden layers or the readout, and $S$ independently chosen locations. Each costs $\log N$ in the scalar designated-coordinate bound and contributes at least $\log(N/[2(S+1)])$ to the embedding entropy. First-layer input coefficients do not participate in this cancellation; their fan-in is only $d+1$.
\end{proof}

If $m\le N<2(S+1)$, the formula itself fits in the prescribed learner. Place its nodes at fixed indices, with the same designated and background laws. There is no embedding mixture, but every designated internal $\log N$ cost is at most $\log(2(S+1))$. Thus the same entropy bound holds.

For a general width-$m$ teacher, let $S$ bound the size of its actual rooted formula expansion. In the large-width case, use that expansion as above. Each unrolled row retains at most $m$ predecessors, so the moment resource continues to use the original width $m$, even though the number of formula occurrences can be much larger. In the small-width case, embed the original shared network instead. Every used internal/readout edge appears in at least one formula occurrence, so there are at most $S$ such edges. For a fully connected teacher the coarse bound is $(H-1)m^2+m\le\sum_{\ell=1}^H m^\ell$. The original implementation has at most $S$ relevant hidden rows and at most $(d+1)S$ first-layer coordinates. The moment argument below applies to both product comparison constructions.

\begin{lemma}[Entropy of the comparison]\label{lem:entropy}
For every finite $N\ge m$, the comparison law satisfies
\begin{equation}
 \KL(Q\|\mu_N)
 \le C(d+1)(S+1)
       \log\frac{C(S+1)(d+1)(B+8)}{\epsilon}.
 \label{eq:comparison-entropy}
\end{equation}
The constant depends only on the fixed L\'evy measure and is independent of depth.
\end{lemma}
\begin{proof}
Each selected background row has at most its fan-in many background coordinates. Their scalar KL bounds therefore sum to at most $C\log(C/\epsilon)$. There are at most $S+1$ selected rows, including the readout. There are at most $S$ designated internal/readout edges and $(d+1)S$ designated first-layer coordinates. Apply their scalar entropy bounds and Lemma~\ref{lem:embedding}, or the small-width bound on $\log N$. Product additivity gives the result for a fixed embedding, followed by~\eqref{eq:embedding-entropy} for the mixture.
\end{proof}

The mixture over embeddings is a comparison measure, not the learner's prior. The prior remains the independent product law $\mu_N$.

\section{A small formula gives a substantial function neighborhood}
\label{sec:formula-uniform-mass}
The fixed-embedding comparison laws are product laws. Only first moments are needed for an expected uniform approximation bound, even for unbounded Lipschitz activations.

\begin{lemma}[Expected uniform comparison error]
\label{lem:formula-uniform-error}
Put $D=C(1+A+b)(m+1)(B+8)$ and $W=C(d+1)D^{H+1}$, with a sufficiently large constant depending only on the fixed L\'evy measure. At comparison precision $0<\epsilon\le1$,
\[
 \E_Q\sup_{x\in[-1,1]^d}|z_w(x)-g(x)|\le W\epsilon.
\]
The constants are independent of depth and learner width.
\end{lemma}
\begin{proof}
Fix an embedding. Define $R_j^1=b+A\sum_{k=0}^d|W^1_{jk}|$ and $R_j^\ell=b+A\sum_k|W^\ell_{jk}|R_k^{\ell-1}$; these bound hidden magnitudes simultaneously on the cube. Let $S_\ell$ be the maximum expected path bound of a selected node and $P_\ell$ the average expected path bound in that layer. Native weights have expected absolute magnitude at most $M_1/F$, designated weights at most $B+8$, and background weights at most $C\epsilon/F$. Independence from preceding layers gives
\begin{align*}
 S_1+P_1&\le C(d+1)(1+A+b)(B+8),\\
 S_\ell&\le b+Am(B+8)S_{\ell-1}+CA P_{\ell-1},\\
 P_\ell&\le S_\ell+b+AM_1P_{\ell-1}.
\end{align*}
The last line separately bounds selected and native rows, using fractions at most one. A sufficiently large fixed coefficient in $D$ gives $S_\ell+P_\ell\le C(d+1)D^\ell$.

Let $e_\ell$ be the maximum expected uniform error of a selected hidden node. The first layer obeys $e_1\le CA(d+1)\epsilon$. Expanding later preactivation differences into teacher coefficients times preceding errors, designated perturbations times new features, and background contributions gives
\[
 e_\ell\le AmB e_{\ell-1}+CA\epsilon[mS_{\ell-1}+P_{\ell-1}].
\]
Here $\E|W-a|\le C\epsilon$ on designated coordinates, and only independence between current weights and preceding layers is used. Taking $D$ to dominate twice the propagation coefficient gives $e_\ell\le C(d+1)D^\ell\epsilon$, with a depth-independent constant. The linear readout adds one factor $D$. This also handles shared-network small-width embeddings. Averaging over embeddings preserves the bound.
\end{proof}

\begin{theorem}[Formula complexity controls negative log prior mass]
\label{thm:uniform-mass}
For a layered formula $g$ with hidden depth $H$, at most $S$ hidden occurrences, maximum hidden-layer width $m$, and coefficients bounded by $B\ge1$, every finite $N\ge m$ and $0<\tau\le1$ satisfy
\[
 -\log\mu_N\{\|z_w-g\|_{\infty,[-1,1]^d}\le\tau\}
 \le C(d+1)(S+1)\log\frac{C(S+1)(d+1)(B+8)W}{\tau},
\]
where $W=C(d+1)[C(1+A+b)(m+1)(B+8)]^{H+1}$.
An arbitrary width-$m$, depth-$H$ network satisfies the same bound with $S=\sum_{\ell=1}^H m^\ell$, at every $N\ge m$.
\end{theorem}
\begin{proof}
Take $\epsilon=\tau/(2W)$, increasing $W$ so this is at most one. Expected uniform approximation and Markov give comparison probability at least one half to the stated event $E$. Binary relative entropy gives
\[
 \KL(Q\|\mu_N)\ge Q(E)\log\frac1{\mu_N(E)}-\log2,
\]
so $-\log\mu_N(E)\le2\KL(Q\|\mu_N)+2\log2$. Apply Lemma~\ref{lem:entropy}, including its shared-network branch.
\end{proof}

\section{Typical prior functions have short formula descriptions}
\label{sec:levy-function-compression}
The small-formula prior-mass theorem has a converse. With high probability, a draw from the same dense prior is uniformly approximable by a finite formula whose size and coefficient precision are bounded independently of learner width. The proof uses the Poisson representation only to describe that draw; it does not modify the prior.

A formula here is an ordered tree of affine summations and the fixed hidden activation $\phi$, with a linear output node. Input leaves are labelled by one of $1,x_1,\ldots,x_d$. A repeated input or hidden computation may occur several times in the tree. Size counts gates and weighted edges. Dyadic coefficient bit length includes the signed numerator and the power-of-two denominator. Zero-input hidden gates are allowed and represent the fixed constant $\phi(0)$.

\begin{theorem}[Uniform approximation of prior draws]
\label{thm:levy-prior-compression}
Fix the L\'evy measure and a globally $A$-Lipschitz activation with $A\ge1$ and $b=|\phi(0)|$. There is a constant $C$, depending only on these fixed quantities, such that the following holds for every input dimension $d$, hidden depth $H\ge1$, finite learner width $N$, and $0<\delta,\tau\le1/2$. Put
\[
 L=H+1+\log\frac1{\delta\tau},\qquad
 S_*=[CL\log(eL)]^{H+2},\qquad B_*=CL^3.
\]
There is a finite dictionary $\mathcal D_{H,d,\delta,\tau}$ of formulas of hidden depth at most $H$, size at most $S_*$, and dyadic coefficient bit length at most $B_*$, independent of $N$, such that
\[
 \mu_N\left\{\inf_{g\in\mathcal D_{H,d,\delta,\tau}}
          \sup_{x\in[-1,1]^d}|z_w(x)-g(x)|\le\tau\right\}
 \ge1-\delta.
\]
Its logarithmic cardinality is at most
\[
 \log|\mathcal D_{H,d,\delta,\tau}|
       \le[1+\log(d+1)][CL\log(eL)]^{H+6},
\]
after increasing $C$. The same conclusion holds after a fixed Lipschitz output map, with adjusted constants and that map appended to each formula.
\end{theorem}

\subsection{Truncating small jumps uniformly on the input cube}
For each weight use its native Poisson representation and remove all jumps with magnitude at most $\epsilon$. Denote the resulting weight by $W^{>\epsilon}$ and the resulting raw network by $z^{>\epsilon}$. The omitted remainder satisfies
\[
 \E|W-W^{>\epsilon}|
 \le\frac1F\int_{|u|\le\epsilon}|u|\mathfrak n(du)
 \le C\epsilon/F,
\]
while both $\E|W|$ and $\E|W^{>\epsilon}|$ are at most $M_1/F$.

For a hidden node define the absolute path bound
\[
 R_j^1=b+A\sum_{k=0}^d|W^1_{jk}|,
 \qquad
 R_j^\ell=b+A\sum_k|W^\ell_{jk}|R_k^{\ell-1}.
\]
It bounds the hidden feature uniformly on the cube. The analogous truncated path bound has the same expected upper bound. Since the total incoming expected absolute weight is at most $M_1$ in every row, layer independence gives
\[
 \E R_j^\ell,\ \E R_j^{\ell,>\epsilon}\le K^\ell,
 \qquad K=1+b+AM_1.
\]
Let $e_\ell$ be the maximal expected uniform difference at a hidden node. Using the full weights on the propagated difference and the omitted remainder on the truncated feature gives
\[
 e_\ell\le AM_1e_{\ell-1}+CA\epsilon K^{\ell-1},
 \qquad e_0=0.
\]
The current layer's weights are independent of all preceding features; independence between its full weight and its own remainder is not needed. Including the linear readout proves
\begin{equation}
 \E\sup_x|z_w(x)-z^{>\epsilon}(x)|
       \le C\epsilon(H+1)K^{H+1}.
 \label{eq:levy-truncation-uniform}
\end{equation}
All constants are independent of width and input dimension. Choose
\[
 \epsilon=\min\left\{\frac12,
          \frac{\delta\tau}{16C(H+1)K^{H+1}}\right\}.
\]
Markov's inequality makes this uniform truncation error at most $\tau/2$ outside a prior event of probability at most $\delta/4$. Also
\begin{equation}
 \lambda:=\Lambda(\epsilon)\le C L.
 \label{eq:levy-row-intensity}
\end{equation}

\subsection{A stopped exploration of the output ancestry}
Across a row of fan-in $F$, the number of retained jumps is exactly Poisson with mean $F\Lambda(\epsilon)/F=\lambda$. These counts are independent between rows. Start at the readout row and reveal only the hidden rows reached by its retained jumps, continuing backwards through the layers. Multiple jumps at one coordinate are kept as separate weighted edges; repeated occurrences of a hidden node reuse its already revealed row. Unrolling this directed acyclic graph produces a formula with the same function.

The choice to inspect a previously unrevealed row depends only on weights in later layers and on already revealed distinct rows. It is therefore independent of that row's incoming jumps. This remains true for the first hidden layer, whose fan-in is $d+1$ instead of $N$.

Choose an integer
\[
 T=\left\lceil C L\log(eL)\right\rceil,
 \qquad S=2(T+1)^{H+1},
\]
with a sufficiently large fixed $C$. Stop exploration at the first row with more than $T$ retained jumps. Before such a failure, no more than $S$ rows can be inspected: a depth-$H+1$ tree with branching at most $T$ has fewer nodes than that bound. For $T\ge e^2\lambda$, the Poisson tail satisfies
$\Pr\{\operatorname{Poisson}(\lambda)>T\}\le e^{-T}$.
A conditional union bound over the stopped exploration therefore gives failure probability at most $S e^{-T}\le\delta/4$. The final inequality follows by increasing the fixed coefficient in $T$: since $H+1\le L$, $\log S=O((H+1)\log(T+1))\le C_1L\log(eL)$, while $\log(1/\delta)\le L$. A fixed sufficiently large coefficient dominates these bounds and the condition $T\ge e^2\lambda$.

The retained jump magnitudes also have a width-independent bound. For $B\ge1$, the intensity of jumps larger than $B$ in one complete row is
\[
 \Lambda(B)\le C B^{-p}.
\]
At most $S$ previously unrevealed rows are examined before the degree stop. The same conditional union argument gives probability at most $CSB^{-p}$ of seeing a jump larger than $B$. Choose
\[
 B=(C S/\delta)^{1/p}+1.
\]
Outside another event of probability at most $\delta/4$, the explored formula has at most $S$ weighted edges and gates, at most $T$ incoming jumps per node, and each edge coefficient has magnitude at most $B$. No statement about all $N H$ rows is needed. In particular there is no hidden $\log N$ factor.

\subsection{Quantization and the finite dictionary}
On the preceding event, round every retained jump coefficient to a dyadic grid of step $q$. Each rounded coefficient has magnitude at most $B+1$. Put
\[
 K_q=1+b+AT(B+1).
\]
The absolute output of a hidden formula of depth $\ell$ is at most $K_q^{\ell}$ on the cube. A telescoping layerwise error estimate therefore bounds the final raw output change by
\[
 q(H+1)T K_q^{H+1}.
\]
Choose a power-of-two grid step satisfying
\[
 q\le\frac{\tau}{4(H+1)T K_q^{H+1}}.
\]
The right side uses the already fixed bound $B+1$, not the actual rounded coefficients, so this is not a circular definition. The quantization error is at most $\tau/4$.

From \eqref{eq:levy-row-intensity}, $T\le CL\log(eL)$ and $S\le[CL\log(eL)]^{H+2}$. Also
\[
 \log(B+1)\le C[\log S+\log(1/\delta)],
\]
so both the numerator and denominator of each rounded coefficient have at most $CL^3$ bits. Indeed the needed grid exponent is at most
\[
 C\left[\log(1/\tau)+(H+1)
                 \{1+\log T+\log(B+1)\}\right]\le CL^3.
\]
Take the dictionary of all formulas obeying these size, depth, coefficient-precision and input-label bounds. A prefix tree description specifies each node's arity, edge coefficients, and any input label. Its total description length is at most
$C S_*[B_*+\log(S_*+1)+\log(d+1)]$.
This gives the stated logarithmic cardinality. The truncation, degree, and magnitude failures have total probability at most $3\delta/4$, and the successful approximation error is at most $3\tau/4$. Enlarging constants gives the theorem with its displayed $\delta,\tau$.

\begin{corollary}[Prior mass bounds approximation complexity]
\label{cor:levy-prior-complexity-converse}
Fix a prescribed input subset $\mathcal X\subseteq[-1,1]^d$, a raw target function $f$ on that set, and $0<\tau\le1/2$. Write $\|\cdot\|_{\infty,\mathcal X}$ for the supremum on this set, and assume the following event is measurable. If
\[
 p=\mu_N\{\|z_w-f\|_{\infty,\mathcal X}\le\tau\}>0,
\]
then $f$ has a uniform $2\tau$ approximation on $\mathcal X$ by a formula of hidden depth at most $H$, size at most
\[
 [C L_p\log(eL_p)]^{H+2},\qquad L_p=H+1+\log\frac1{p\tau},
\]
and dyadic coefficient bit length at most
\[
 C\left(H+1+\log\frac1{p\tau}\right)^3.
\]
The bounds are independent of learner width. Combined with Theorem~\ref{thm:uniform-mass}, this gives a two-sided approximate function-prior relation at fixed depth: a short bounded-coefficient formula has a prior neighborhood with controlled negative log mass, and a neighborhood of appreciable prior mass contains a function approximable by a short finite-precision formula.
\end{corollary}
\begin{proof}
Apply Theorem~\ref{thm:levy-prior-compression} with failure probability $p/2$ and approximation tolerance $\tau$. Its successful set intersects the stated neighborhood, because that neighborhood has mass $p$ and the exceptional set has mass at most $p/2$. For a network in the intersection, the triangle inequality gives a $2\tau$ approximation to $f$. Substitute $\delta=p/2$ into the resource bounds and absorb the factor two into $C$.
\end{proof}

For ReLU, this also gives an ordinary Boolean-circuit consequence on the Boolean cube. If the target takes values in $\{-1,1\}$ and $2\tau<1$, the rational ReLU formula in the corollary has exactly the target's sign at every Boolean input. At each gate, rational fixed-coefficient multiplication, addition and comparison with zero can be implemented by Boolean arithmetic circuits. With dyadic weights of $B_*$ bits, all intermediate integer bit lengths are polynomial in $H$, $B_*$ and $\log(S_*+1)$: a common denominator at layer $\ell$ has at most $\ell B_*$ bits, and the coefficient magnitude and fan-in bounds control numerator growth. Schoolbook arithmetic consequently gives a Boolean circuit of size polynomial in $d,S_*,H,B_*$. Thus, at fixed depth and fixed margin tolerance, a Boolean function with appreciable raw-function prior mass has Boolean circuit complexity polynomial in its negative log prior mass and input dimension. This statement concerns approximation neighborhoods of real outputs; it does not identify their probability with the mass of an unrestricted sign event lacking a positive output margin.

\section{An exact Boolean sign prior under an all-input margin condition}
\label{sec:levy-margin-converse}
The compression theorem also applies directly to a margin sign event. It does not require the raw output to stay near a prescribed amplitude or to have an upper bound.

\begin{theorem}[Circuit complexity from margin-sign prior mass]
\label{thm:margin-sign-converse}
Use ReLU and a prescribed Boolean input set $\mathcal X\subseteq\{-1,1\}^d$. Let $f:\mathcal X\to\{-1,1\}$ and $\kappa>0$. If
\[
 p=\mu_N\{f(x)z_w(x)\ge\kappa\text{ for every }x\in\mathcal X\}>0,
\]
put $\tau=\min\{1/2,\kappa/2\}$ and $L_p=H+1+\log(1/(p\tau))$. Then a rational ReLU formula of size at most
$[CL_p\log(eL_p)]^{H+2}$ and coefficient bit length at most $CL_p^3$ has sign exactly $f$ on $\mathcal X$. Consequently $f$ has a Boolean circuit of size polynomial in $d$, that formula size, $H$, and the coefficient bit length.

The same conclusion holds for the prior conditioned on the all-input margin event
\[
 M_\kappa=\{|z_w(x)|\ge\kappa\text{ for every }x\in\mathcal X\},
\]
if $\log(1/p)$ in $L_p$ is replaced by $\log(1/p)+Q_0$. Here $p$ is the conditional probability of sign function $f$, and a bound independent of learner width is
\begin{align*}
 Q_0\le C(d+1)(H+1)\big[&H+\log(d+H+2)\\
 &+(H+2)\log(1+\kappa)+\log(1+\kappa^{-1})\big].
\end{align*}
At fixed hidden depth and margin, both conclusions are polynomial circuit-complexity bounds in input dimension and negative log function-prior mass.
\end{theorem}
\begin{proof}
Apply Theorem~\ref{thm:levy-prior-compression} with failure probability $p/2$ and approximation error $\tau$. Its good set intersects the margin sign event. At a network in that intersection, the approximating formula $g$ satisfies
\[
 f(x)g(x)\ge\kappa-\tau\ge\kappa/2>0.
\]
This proves exact sign agreement. The rational ReLU-to-Boolean conversion was detailed after Corollary~\ref{cor:levy-prior-complexity-converse}; it does not require values of $g$ to be bounded above.

For the conditioned prior, a width-one ReLU chain computes the constant $2\kappa$: its first-layer bias is one, all later hidden weights are one, and its readout coefficient is $2\kappa$. It has $H$ hidden occurrences and coefficient bound $B_0=\max\{1,2\kappa\}$. A uniform raw neighborhood of radius $t_0=\min\{1,\kappa/2\}$ around this function lies in $M_\kappa$, even on the whole input cube. Theorem~\ref{thm:uniform-mass} gives
\[
 \mu_N(M_\kappa)\ge e^{-Q_0}
\]
with the displayed bound after expanding its constants and logarithms. In particular this event has positive probability at every finite width. If its conditional sign mass is $p$, the original joint margin sign event has mass
$p\mu_N(M_\kappa)\ge p e^{-Q_0}$. Apply the first part with this lower bound.
\end{proof}

Conversely, a ReLU formula implementing a Boolean target with margin at least $2\kappa$ has a uniform neighborhood contained in its correct-sign margin event. Theorem~\ref{thm:uniform-mass} bounds the negative log probability of that event. Conditioning on $M_\kappa$ can only increase its probability. Thus the all-input-margin convention also has two directions: small margin formulas receive controlled prior mass, and sign functions with appreciable conditional prior mass have small Boolean circuits at fixed depth. An unrestricted sign event without a positive margin is not covered by this converse.

\section{An explicit-density alternative: signed beta-prime weights}
\label{sec:beta-prime}
The comparison and compression arguments use scalar moment, small-ball and tail estimates rather than infinite divisibility itself. A simple explicit density satisfies the same estimates.

Fix $p>2$ and $0<c_0\le1$. For fan-in $F$, put $a_F=c_0/F$ and give each coordinate the independent symmetric density
\begin{equation}
 \varrho_F(w)=
 \frac{|w|^{a_F-1}(1+|w|)^{-a_F-p}}
      {2\,\mathrm B(a_F,p)},\qquad w\ne0,
 \label{eq:beta-prime-prior}
\end{equation}
where $\mathrm B$ is Euler's beta function. Equivalently, $|W|$ has the beta-prime law with shape parameters $(a_F,p)$ and its sign is an independent fair coin. This is an absolutely continuous full-support distribution with no zero atom. Its density is sharply concentrated near zero and has a fixed polynomial tail. Fan-in changes its first shape parameter, not the architecture.

\begin{proposition}[The same function-prior guarantees for an explicit density]
\label{prop:beta-prime}
All uniform-neighborhood, native compression, complexity-converse and margin-conditioned sign-prior conclusions of this paper remain valid, with modified fixed constants, if every L\'evy increment coordinate is replaced by~\eqref{eq:beta-prime-prior}.
\end{proposition}
\begin{proof}
We verify the scalar interface. Write $a=a_F\in(0,1]$ and
$\mathcal B=\mathrm B(a,p)$. The beta integral gives
\[
 \E|W|=\frac{a}{p-1},\qquad
 \E W^2=\frac{a(a+1)}{(p-1)(p-2)}.
\]
Thus the native first absolute and second moments are at most $C/F$.
Also $c_p/a\le\mathcal B\le1/a$, with fixed $c_p>0$. For the lower bound integrate the beta-prime density kernel on $(0,1)$; for the upper bound use the beta integral on $(0,1)$ with factor $(1-u)^{p-1}\le1$.

For $0<\epsilon\le1$, let $P_\epsilon=\Pr(|W|\le\epsilon)$.
The integral above $\epsilon$ gives
\[
 1-P_\epsilon\le Ca\{1+\log(1/\epsilon)\}.
\]
Indeed the integral on $(\epsilon,1)$ is at most
$\log(1/\epsilon)$, the integral on $(1,\infty)$ is at most $1/p$,
and $\mathcal B^{-1}\le Ca$.
A second bound, obtained by integrating below $\epsilon$, is
\[
 P_\epsilon\ge
 \frac{2^{-p-1}\epsilon^a}{a\mathcal B}
 \ge2^{-p-1}\epsilon^a.
\]
When $Ca\{1+\log(1/\epsilon)\}\le1/2$, the first bound and
$-\log(1-u)\le2u$ give the desired logarithmic conditioning cost.
Otherwise the second bound gives
$-\log P_\epsilon\le a\log(1/\epsilon)+(p+1)\log2$;
the fixed additive constant is absorbed by
$a\{1+\log(1/\epsilon)\}$ in this case. Therefore
\begin{equation}
 -\log P_\epsilon\le\frac C F\log(C/\epsilon).
 \label{eq:beta-prime-small-ball}
\end{equation}

Condition the coordinate directly on $|W|\le\epsilon$.
Symmetry preserves mean zero. For $r=1,2$, the ratio of its truncated moment integral to the lower bound on $P_\epsilon$ gives
\[
 \E[|W|^r\mid |W|\le\epsilon]
 \le 2^{p+1}\frac{a}{a+r}\epsilon^r
 \le C\epsilon^r/F.
\]
Together with~\eqref{eq:beta-prime-small-ball}, these are exactly the background support, entropy and moment estimates used in the network proof.

For a designated coefficient $u$, $|u|\le B$, use the unbiased separated interval law from Lemma~\ref{lem:levy-designated}, without adding a small-jump remainder. Its support lies in
$\{\epsilon<|w|\le B+8\}$, its density is at most $1/\epsilon$,
and its variance around $u$ is at most $C\epsilon^2$.
On this support,
\[
 \varrho_F(w)\ge
 \frac{a}{2(B+8)(B+9)^{p+1}}.
\]
Hence its relative entropy to $\varrho_F$ is at most
$\log F+C\log(C(B+8)/\epsilon)$.
All designated-support and unbiasedness properties are unchanged.

The comparison embedding and uniform-neighborhood proofs use only these scalar estimates and therefore apply unchanged.

For native compression, truncate whole coordinates of magnitude at most $\epsilon$, rather than deleting individual Poisson jumps. The density bounds imply
\[
 \E[|W|\mathbf1_{|W|\le\epsilon}]\le C\epsilon/F,
 \quad F\Pr(|W|>\epsilon)\le C\log(C/\epsilon),
 \quad F\Pr(|W|>B)\le CB^{-p}\quad(B\ge1).
\]
The expected uniform truncation error is consequently the same as in Theorem~\ref{thm:levy-prior-compression}. A fresh row has a binomial number of retained coordinates with mean at most $C\log(C/\epsilon)$. Its exponential moment is bounded by that of a Poisson variable of the same mean. Hence the stopped-ancestry degree and magnitude bounds, and the ensuing quantization and dictionary arguments, remain valid with the same resource exponents. Fresh-row independence is unchanged. The neighborhood and margin converses depend only on the resulting high-probability compression and the already verified neighborhood lower bound. No Poisson representation is needed for this alternative.
\end{proof}

The density behaves near zero like $|w|^{-1+c_0/F}$, with integrable central concentration and no exact sparsity atoms. It can be viewed as continuous shrinkage whose strength is tied to fan-in. This is a mathematical prior construction, not an empirical claim about trained-network weights.

Its wide-row limit has the same L\'evy mechanism. For fixed $w\ne0$,
\[
 F\varrho_F(w)\longrightarrow
 \frac{c_0}{2|w|(1+|w|)^p}.
\]
Dominated convergence against $\cos(tw)-1$, using the quadratic bound near zero and the integrable tail bound, shows that a sum of $F$ such independent coordinates converges to the non-Gaussian infinitely divisible law with this L\'evy measure. Its variance remains of order one. Exact finite-width infinite divisibility is special to increment laws; none of the conclusions transferred above requires it.

\section{The wide limit is a random formula}
\label{sec:wide-formula-limit}
The distinction between shared circuits and formulas has a structural explanation for independent sparse coordinates. At any fixed resolution, the ancestors of the output become a tree as learner width tends to infinity. This does not by itself give a lower bound on learning a particular function: a function may have several very different implementations.

\begin{theorem}[A random-formula limit]
\label{thm:wide-formula-limit}
Fix input dimension $d$, hidden depth $H$, the globally Lipschitz activation, and the scalar L\'evy prior. The pushforwards of $\mu_N$ to raw network functions converge weakly on $C([-1,1]^d)$ as $N\to\infty$. For each fixed $\epsilon>0$, their approximation obtained by deleting internal/readout jumps of magnitude at most $\epsilon$ converges in total variation to a finite random formula law. The limiting formula has Poisson offspring of mean
\[
 \Lambda_\epsilon=\mathfrak n\{|u|>\epsilon\}
\]
at each internal row and at the readout, independent edge weights with law
$\mathfrak n(du)\mathbf1_{|u|>\epsilon}/\Lambda_\epsilon$, and independent native first-layer random features at its leaves. Its hidden depth is at most $H$. For the truncated laws, a total-variation error bound is
\[
 \frac{C(H+1)^3(1+\Lambda_\epsilon)^{2H}}{N},
\]
truncated at one. The full limiting prior is the limit in distribution of these finite random-formula laws as $\epsilon\downarrow0$. For the bounded metric $d(f,g)=\min\{1,\|f-g\|_\infty\}$, its Wasserstein distance from the width-$N$ function prior is at most $C_H(1+\log N)^{2H}/N$ for $N\ge2$, with fixed-depth constants independent of width.
\end{theorem}
\begin{proof}
Keep the first layer native, and truncate only the remaining hidden layers and readout. A truncated incoming row is exactly a Poisson collection of jumps of total intensity $\Lambda_\epsilon$: each jump chooses its child index uniformly from $N$ coordinates, and its magnitude and sign have the displayed law. Incoming rows are independent.

Generate instead an independent depth-$H$ Poisson tree, with fresh first-layer random features at the bottom. Couple it to the finite-width network by assigning each newly generated hidden vertex an independent uniform index in its layer. Until an index is reused in the same layer, all visited rows are fresh and their laws agree exactly with the corresponding independent tree rows. Conditional on the tree size $S$, a union bound gives collision probability at most $S^2/(2N)$. For a Poisson branching process, the generation second moments satisfy
\[
 \E Z_{j+1}^2=\Lambda_\epsilon^2\E Z_j^2
                         +\Lambda_\epsilon\E Z_j,
 \qquad \E Z_j=\Lambda_\epsilon^j.
\]
Consequently $\E S^2\le C(H+1)^3(1+\Lambda_\epsilon)^{2H}$. On the event of no collisions the two functions agree identically, which proves the total-variation claim. Multiple jumps choosing one coordinate are included among the collision events.

The first-moment truncation argument in the compression theorem also applies when the first layer is left native. Uniformly in $N$,
\[
 \E\|z_N-z_N^{(\epsilon)}\|_\infty
 \le C\epsilon(H+1)K^{H+1}
\]
for a fixed $K$ depending only on the activation and scalar law. Use the bounded complete metric
$d(f,g)=\min\{1,\|f-g\|_\infty\}$ on the function space. The corresponding Wasserstein distance between the full and truncated laws is bounded by the same expression. The limiting tree laws are therefore Cauchy in this metric as $\epsilon\downarrow0$: compare two truncation levels through the finite-width coupling and then let $N\to\infty$. Completeness of the Wasserstein space over this bounded complete Polish metric gives their limiting law. The same triangle inequality proves convergence of the original function priors to it. In fact the distance from a truncated tree law to the limiting law is at most $C\epsilon(H+1)K^{H+1}$. Taking $\epsilon=1/N$ and using $\Lambda_\epsilon\le C\log(C/\epsilon)$ gives the stated quantitative rate.
\end{proof}

\begin{corollary}[Fixed-data posterior limit]
Append any continuous final link into $[-1,1]$, use a fixed Gaussian likelihood variance, and fix a finite observed data set. The corresponding posterior laws of raw functions converge weakly to the posterior obtained by tilting the limiting prior in Theorem~\ref{thm:wide-formula-limit}. Bounded posterior-mean predictions converge at every input.
\end{corollary}
\begin{proof}
For the fixed data set, the likelihood is a bounded, strictly positive continuous functional of the raw function. Its normalizer converges to a positive number. Integrating bounded continuous functions against the normalized likelihood proves posterior convergence; evaluating the bounded final link at an input gives the prediction assertion.
\end{proof}

\begin{proposition}[The finite-width cost of a shared motif]
\label{prop:shared-motif}
Fix a layered output-relevant graph $G$ with $V$ distinct hidden vertices and $E$ internal/readout edges. Put $J=E-V\ge0$. At resolution $0<\epsilon\le1$, let $A_G$ be the event that the network contains an injective embedding of $G$ in which every indicated weight has absolute value greater than $\epsilon$. Under either independent scalar prior in this paper,
\[
 \mu_N(A_G)\le\lambda_\epsilon^E N^{-J},
 \qquad\lambda_\epsilon=C\log(C/\epsilon).
\]
Additional coefficient restrictions only decrease this probability. With a final link into $[-1,1]$ and Gaussian likelihood written as $\exp[-\sum_i(f(X_i)-Y_i)^2/T]$, its posterior probability obeys
\[
 \Pi_{S_n}(A_G)
 \le \exp\left[\frac4T\sum_{i=1}^n(1+|Y_i|)\right]
             \lambda_\epsilon^E N^{-J}.
\]
In particular, for labels bounded by $B_Y$, assigning posterior mass at least $\delta$ to this shared motif requires
\[
 n\ge\frac{T}{4(1+B_Y)}
   \left[J\log N-E\log\lambda_\epsilon-\log(1/\delta)\right]_+.
\]
\end{proposition}
\begin{proof}
For signed beta-prime weights the scalar estimates above give
$\Pr(|W|>\epsilon)\le\lambda_\epsilon/N$. For L\'evy increments, split at jump size $\epsilon/2$. The probability of a larger jump is at most $\Lambda(\epsilon/2)/N$. The remaining absolute jump sum has expectation at most $C\epsilon/N$, so Markov's inequality bounds the probability it exceeds $\epsilon/2$ by $C/N$. This proves the same estimate.

There are at most $N^V$ injective placements of the hidden vertices. For each placement the $E$ distinct indicated coordinates are independent, so their probability is at most $(\lambda_\epsilon/N)^E$. The union bound proves the first claim. The squared-loss oscillation over predictions in $[-1,1]$ is at most $4(1+|Y_i|)$ for each observation. The posterior-to-prior likelihood ratio is therefore at most the displayed exponential factor, giving the remaining conclusions.
\end{proof}

The factor $N^{-J}$ matches the uncancelled $J\log N$ in a direct shared-circuit comparison. A tree embedding has one outgoing edge per hidden occurrence and $J=0$. A shared computation has at least one extra outgoing edge and $J>0$. Independent coordinates make each such edge rare separately, while a vertex contributes only one choice of location.

This is an obstruction to retaining a prescribed shared implementation at fixed resolution, not a proof that its function needs $\Omega(\log N)$ samples. A formula can compute the same function, and the posterior may use that representation instead. Proving a function-level lower bound would require an additional lower bound on all sufficiently accurate formula representations and their prior mass. The limiting tree description identifies that remaining question without assuming such a complexity separation.


\section{Interpretation and verification}
The central concentration and rare order-one jumps have different roles. Background suppression costs only a logarithm of precision. A designated coefficient costs a factor of order $1/N$, but every hidden formula occurrence has $N$ possible locations; the comparison entropy accounts for those locations exactly. Conversely, a typical output depends, at any fixed resolution, on a finite Poisson ancestry tree even when the surrounding architecture is arbitrarily wide.

This is a continuous shrinkage prior. It is not a small Gaussian perturbation: rare order-one jumps persist in wide sums, although coordinate variance is proportional to inverse fan-in. The sum of $F$ independent fan-in coordinates has exactly the non-Gaussian law $L_1$ for every $F$. An order-one initialization variance alone therefore does not imply a Gaussian-process limit.

The proofs have undergone independent AI-agent mathematical review and have not been formalized in Lean. The verification record identifies the exact reviewed sources. These are prior-probability and approximation-complexity statements, not efficient posterior-sampling guarantees.
\end{document}
